% Evidara V19.1 — one year / one paper / one set per file (or ZIP)
% Put any referenced images in the same ZIP, for example images/q001.png.

\paper{NEET 2026 Main Paper}
\exam{NEET}
\year{2026}
\variant{Main}
\code{NEET-2026-MAIN}
\duration{180}
\grade{Grade 11-12}
\pyq{true}
\sourcekey{NEET_2026_MAIN}
\description{NEET 2026 previous-year question paper.}
\instructions{Attempt the paper according to the original examination pattern.}

\begin{question}
\id{NEET-2026-Q001}
\number{1}
\section{Physics}
\subject{Physics}
\chapter{Units and Measurements}
\topic{Dimensional Analysis}
\questiontype{single_correct}
\difficulty{moderate}
\stem{The dimensions of Planck's constant are the same as the dimensions of $X$. Identify $X$.}
\option{A}{Force}
\option{B}{Angular momentum}
\option{C}{Power}
\option{D}{Pressure}
\answer{B}
\solution{Planck's constant has dimensions $ML^2T^{-1}$, which are the dimensions of angular momentum.}
\marks{4}{1}
\end{question}

\begin{question}
\id{NEET-2026-Q002}
\number{2}
\section{Physics}
\subject{Physics}
\chapter{Motion in a Straight Line}
\topic{Kinematics}
\questiontype{single_correct}
\difficulty{moderate}
\stem{A particle starts with velocity $u$ and moves with uniform acceleration $a$ for time $t$. Its final velocity is}
\option{A}{$v=u-at$}
\option{B}{$v=u+at$}
\option{C}{$v=ut+a$}
\option{D}{$v=u+a/t$}
\answer{B}
\solution{For constant acceleration, \(a=\frac{v-u}{t}\), therefore \(v=u+at\).}
\marks{4}{1}
\end{question}

\begin{question}
\id{NEET-2026-Q003}
\number{3}
\section{Chemistry}
\subject{Chemistry}
\chapter{Chemical Bonding}
\topic{Molecular Geometry}
\questiontype{single_correct}
\difficulty{moderate}
\stem{Use the molecular diagram below and choose the correct geometry. \includegraphics{images/q003.png}}
\option{A}{Linear}
\option{B}{Trigonal planar}
\option{C}{Tetrahedral}
\option{D}{Octahedral}
\answer{C}
\solution{Count the electron domains around the central atom and apply VSEPR theory.}
\marks{4}{1}
\end{question}

% If a question is entirely mathematical, use an explicit LaTeX block:
% \begin{question}
% \number{4}
% \section{Physics}
% \subject{Physics}
% \questiontype{numerical}
% \latex{\frac{1}{R}=\frac{1}{R_1}+\frac{1}{R_2}}
% \numericalanswer{5}
% \solution{Use the parallel-resistance relation.}
% \marks{4}{0}
% \end{question}
